\providecommand{\Eexp}{\mathcal E_{\mathrm{exp}}}

\section{Spectral stability, arithmetic Clark data, and a second
recovery of \texorpdfstring{$\gamma$}{gamma}}
\label{sec:ssc-spectral-stability-clark}

The Joukowski quotient has a positive Laplace representation which places
its zero divisor strictly in the left half-plane.  After rotation and a
centering exponential, this stability becomes a Hermite--Biehler
inequality.  The resulting meromorphic inner function has two complementary
features: its values at algebraic real points are algebraically independent,
whereas their large-point phase recovers Euler's constant.  We develop these
statements in a form which keeps the unconditional arithmetic assertions
separate from the still open arithmetic of \(\gamma\).

\subsection{A positive Bessel--Laplace representation}

\begin{proposition}[Bessel--Laplace representation]
\label{prop:ssc-bessel-laplace}
For every \(w\in\C\),
\begin{equation}\label{eq:ssc-bessel-laplace}
 H'(w)=\int_0^1 e^{-sw}
 J_0\!\left(2\sqrt{s(1-s)}\right)\,ds.
\end{equation}
The kernel in \eqref{eq:ssc-bessel-laplace} is strictly positive on
\([0,1]\).
\end{proposition}

\begin{proof}
Set \(t=2s-1\).  The classical Sonine integral, obtained for example by
expanding the exponential and the Bessel function and integrating
termwise, is
\[
 \int_{-1}^{1}e^{at}J_0\!\left(b\sqrt{1-t^2}\right)\,dt
 =2\frac{\sinh\sqrt{a^2-b^2}}{\sqrt{a^2-b^2}},
\]
with removable continuation when \(a^2=b^2\).  Taking
\(a=-w/2\) and \(b=1\) gives
\[
 \int_0^1 e^{-sw}J_0\!\left(2\sqrt{s(1-s)}\right)\,ds
 =2e^{-w/2}
   \frac{\sinh(\frac12\sqrt{w^2-4})}{\sqrt{w^2-4}},
\]
which is \eqref{eq:Hprime-sinh}.  Finally
\(0\leq2\sqrt{s(1-s)}\leq1\), and the alternating series for \(J_0\)
shows that \(J_0(x)>0\) for \(0\leq x\leq1\).
\end{proof}

\begin{theorem}[Rightmost point of every real fibre]
\label{thm:ssc-rightmost-fibre}
Let \(a\in\mathbb R\).  If
\[
 H(w)=H(a),\qquad w\ne a,
\]
then
\begin{equation}\label{eq:ssc-rightmost-fibre}
 \Re w<a.
\end{equation}
Thus \(a\) is the unique rightmost point of its full complex fibre.
\end{theorem}

\begin{proof}
Integrating \eqref{eq:ssc-bessel-laplace} from \(a\) to \(w\) gives
\begin{equation}\label{eq:ssc-fibre-integral}
 H(w)-H(a)=\int_0^1 e^{-as}
 \frac{1-e^{-s(w-a)}}{s}
 J_0\!\left(2\sqrt{s(1-s)}\right)\,ds,
\end{equation}
where the integrand is continued at \(s=0\).  Put \(d=w-a\).  If
\(\Re d\geq0\), then
\[
 \Re\bigl(1-e^{-sd}\bigr)
 =1-e^{-s\Re d}\cos(s\Im d)
 \geq1-e^{-s\Re d}\geq0.
\]
All remaining factors in \eqref{eq:ssc-fibre-integral} are positive.
If the integral vanishes, its real part therefore vanishes identically.
This first forces \(\Re d=0\), and then
\(1-\cos(s\Im d)=0\) for every \(s\in[0,1]\), which forces
\(\Im d=0\).  Hence \(d=0\), contrary to the hypothesis.  Thus
\(\Re d<0\).
\end{proof}

\begin{corollary}[Hurwitz stability of the level-zero divisors]
\label{cor:ssc-hurwitz-stability}
Let \(\omega_0\) denote the unique real zero of \(H\).  Then
\(\omega_0<0\), and every other zero \(\rho\) of \(H\) satisfies
\[
 \Re\rho<\omega_0<0.
\]
All zeros of \(H\) are simple and transcendental.  Moreover every zero
\(z\) of
\[
 \Psi(z)=H(z+z^{-1})
\]
is simple, transcendental, and lies in the open left half-plane.
\end{corollary}

\begin{proof}
The real diffeomorphism in Corollary~\ref{cor:H-critical-lattice}, together
with \(H(0)=2\Cin(1)>0\), gives the unique real zero \(\omega_0<0\).
Theorem~\ref{thm:ssc-rightmost-fibre} gives the asserted strict inequality
for every other zero.  The critical points of \(H\) are the purely
imaginary points \eqref{eq:H-critical-points}, so none is a zero; the real
zero is simple because \(H'(\omega_0)>0\).  Transcendence follows from
Corollary~\ref{cor:H-algebraic-levels} at the algebraic level zero.

If \(\Psi(z)=0\), put \(\rho=J(z)\).  Since
\[
 \Re J(z)=\Re z\,(1+|z|^{-2}),
\]
the inequality \(\Re\rho<0\) implies \(\Re z<0\).  The points
\(z=\pm1\), at which \(J'\) vanishes, are not zeros because
\(H(2)>0\) and \(H(-2)<0\).  Hence simplicity follows from the chain
rule.  Finally an algebraic zero \(z\) would make \(J(z)\) an algebraic
zero of \(H\), which has just been excluded.
\end{proof}

\subsection{The Euler point as an explicit
\texorpdfstring{$E$--$D$}{E--D} collision}

Recall that \(w_*\in(-2,2)\) is the unique real solution of
\(H(w)=2\gamma\), and let \(\alpha_*\) and \(\alpha_*^{-1}\) be its two
preimages under \(J\).  Define
\begin{align}
 \mathcal G(u)
 &:=eH'(2+u)
   =e^{-u/2}\sum_{m\geq0}
     \frac{\bigl(u(u+4)/4\bigr)^m}{(2m+1)!},
     \label{eq:ssc-G-definition}\\
 \mathcal K(u)
 &:=-\frac12\int_0^u\mathcal G(t)\,dt.
     \label{eq:ssc-K-definition}
\end{align}
Both functions are entire, $D$-finite, and have rational Taylor
coefficients.  Direct integration gives the useful identity
\begin{equation}\label{eq:ssc-K-H-identity}
 \mathcal K(u)=\frac e2\bigl(H(2)-H(2+u)\bigr).
\end{equation}

There is a tempting but false shortcut here: rational Taylor coefficients
and holonomicity do not make \(\mathcal K\) a Siegel \(E\)-function.

\begin{proposition}[The auxiliary primitive is not an \(E\)-function]
\label{prop:ssc-K-not-E}
Neither \(\mathcal G\) nor \(\mathcal K\) is a Siegel \(E\)-function.
Nevertheless, if \(u\in\Qbar\), then \(\mathcal K(u)\) belongs to the ring
\(\mathbf E\) of finite \(E\)-values.
\end{proposition}

\begin{proof}
Writing \(y=\mathcal G\), the differential equation
\eqref{eq:H-third-order-ode} shifted by \(w=2+u\) gives
\begin{equation}\label{eq:ssc-G-minimal-operator}
 u(u+2)(u+4)y''
 +(u^3+8u^2+16u+12)y'
 +(u^2+3u+4)y=0.
\end{equation}
The function \(y=eH'(2+u)\) has infinitely many zeros, by
\eqref{eq:H-critical-points}.  It therefore cannot satisfy a first-order
homogeneous equation over \(\Qbar(u)\): a nonzero entire solution of such an
equation can have zeros only among the finitely many poles of its rational
logarithmic derivative.  Thus the order-two operator in
\eqref{eq:ssc-G-minimal-operator} is minimal.  At \(u=-4\) its indicial
polynomial is
\[
 8r(r-1)+12r=4r(2r+1),
\]
so its local exponents are \(0\) and \(-1/2\).  Hence \(-4\) is a
non-apparent nonzero finite singularity.  The finite-singularity theorem for
minimal annihilators of \(E\)-functions says that every nonzero finite
singularity is apparent \cite{Andre2000}.  Consequently
\(\mathcal G\) is not an \(E\)-function.  Since \(E\)-functions are closed
under differentiation and \(\mathcal K'=-\mathcal G/2\), neither is
\(\mathcal K\).

For the last assertion, choose algebraic \(a\ne0\) satisfying
\(a+a^{-1}=2+u\).  Equations \eqref{eq:H-descent} and
\eqref{eq:ssc-K-H-identity} give
\begin{equation}\label{eq:ssc-K-finite-E-values}
 \mathcal K(u)=\frac e2
 \bigl(2\Ein(1)-\Ein(a)-\Ein(a^{-1})\bigr).
\end{equation}
The right side belongs to \(\mathbf E\), since \(e\) and the three displayed
finite algebraic-point values are \(E\)-values and \(\mathbf E\) is a ring.
\end{proof}

\begin{proposition}[Equivalent algebraic witnesses]
\label{prop:ssc-euler-equivalences}
The following assertions are equivalent:
\begin{enumerate}[label=\textup{(\roman*)}]
 \item \(w_*\in\Qbar\);
 \item \(\alpha_*\in\Qbar\);
 \item there exists
       \(a\in\Qbar^\times\setminus(-\infty,0]\) such that, on the
       principal branch,
       \[
        E_1(a)+E_1(a^{-1})=0;
       \]
 \item there exists \(u\in\Qbar\cap(-4,0)\) such that
       \[
        \mathcal K(u)=\delta.
       \]
\end{enumerate}
If these conditions hold, then \(\gamma\) and \(\delta\) are
algebraically independent.
\end{proposition}

\begin{proof}
The equivalence of (i) and (ii) follows from
\(\alpha_*^2-w_*\alpha_*+1=0\) and
\(w_*=\alpha_*+\alpha_*^{-1}\).  For a point in the domain specified in
(iii), the inverse-compatible principal branches satisfy
\[
 E_1(z)+E_1(z^{-1})=H(J(z))-2\gamma
\]
and this identity shows that (ii) implies (iii).  Conversely an algebraic
solution in (iii)
gives an algebraic zero \(J(a)\) of \(H-2\gamma\).  By
Theorem~\ref{thm:H-euler-level}, every algebraic zero of this Euler-level
translate must equal \(w_*\), proving (i); the same theorem then forces
\(a\in\{\alpha_*,\alpha_*^{-1}\}\).

Put \(u_*=w_*-2\in(-4,0)\).  Since
\[
 H(2)=2\Ein(1)=2\gamma+\frac{2\delta}{e},
\]
identity \eqref{eq:ssc-K-H-identity} gives
\(\mathcal K(u_*)=\delta\), so (i) implies (iv).  Conversely, if
\(u\in(-4,0)\) satisfies \(\mathcal K(u)=\delta\), the same identity
gives \(H(2+u)=2\gamma\).  Uniqueness on the real axis forces
\(2+u=w_*\), proving (i).  The final assertion is the witness
amplification statement in Theorem~\ref{thm:H-euler-level}.
\end{proof}

\begin{corollary}[Exact \(E\)--\(D\) consequence and the Euler dichotomy]
\label{cor:ssc-ED-dichotomy}
Let \(\mathbf E\) and \(\mathbf D\) denote the rings of \(E\)-values and
arithmetic-Gevrey values of order one, respectively, in the sense of
Fischler--Rivoal \cite{FR2024ED}.  If \(w_*\) is algebraic, then
\[
 \delta\in(\mathbf E\cap\mathbf D)\setminus\Qbar.
\]
Consequently the conjectural equality
\(\mathbf E\cap\mathbf D=\Qbar\) implies that \(w_*\) is
transcendental.  Unconditionally, exactly one of the following alternatives
holds:
\begin{enumerate}[label=\textup{(\alph*)}]
 \item every zero of \(H-2\gamma\), and every zero of the reciprocal
       Euler equation, is transcendental;
 \item the Euler level has the unique algebraic reciprocal pair
       \(\{\alpha_*,\alpha_*^{-1}\}\), and \(\gamma,\delta\) are
       algebraically independent.
\end{enumerate}
\end{corollary}

\begin{proof}
Under the algebraicity hypothesis,
Proposition~\ref{prop:ssc-euler-equivalences} writes \(\delta\) as
\(\mathcal K(u_*)\) at an algebraic point.  Formula
\eqref{eq:ssc-K-finite-E-values}, rather than the false assertion that
\(\mathcal K\) itself is an \(E\)-function, gives
\(\delta\in\mathbf E\).  The inclusion
\(\delta\in\mathbf D\) is the standard Euler-series realization of the
Gompertz constant, and the algebraic independence conclusion in
Proposition~\ref{prop:ssc-euler-equivalences} shows that \(\delta\) is
transcendental.  This proves the first assertion.

Theorem~\ref{thm:H-euler-level} says that \(w_*\) is the only possible
algebraic zero at the Euler level.  If it is transcendental, an algebraic
reciprocal zero is impossible because its Joukowski image would be
algebraic.  If it is algebraic, its two Joukowski preimages are algebraic,
they are the only algebraic reciprocal zeros, and witness amplification
gives the stated algebraic independence.
\end{proof}

\subsection{A level-zero de Branges--Clark model}

For an entire function \(E\), write
\(E^\#(z)=\overline{E(\bar z)}\).  Define
\begin{equation}\label{eq:ssc-E0-Theta}
 E_0(z)=e^{-iz/2}H(-iz),\qquad
 \Theta(z)=\frac{E_0^\#(z)}{E_0(z)}
          =e^{iz}\frac{H(iz)}{H(-iz)}.
\end{equation}

\begin{theorem}[Hermite--Biehler and Clark structure]
\label{thm:ssc-HB-Clark}
The function \(E_0\) is a Cartwright-class Hermite--Biehler function of
exponential type \(1/2\).  Consequently \(\Theta\) is a meromorphic inner
function in the upper half-plane and has zero mean type.  The zeros of
\(\Theta\) are the points \(-i\rho\), where \(\rho\) ranges over the zeros
of \(H\).
Thus \(E_0\) defines a de Branges space, and for each
\(\alpha\in\mathbb T\) the real solutions of
\begin{equation}\label{eq:ssc-Clark-level}
 \Theta(x)=\alpha
\end{equation}
form the corresponding discrete Clark, equivalently self-adjoint, spectrum.
\end{theorem}

\begin{proof}
Corollary~\ref{cor:ssc-hurwitz-stability} places every zero \(\rho\) of
\(H\) in the open left half-plane.  Hence the zeros \(i\rho\) of \(E_0\)
lie in the lower half-plane, and \(E_0\) has no real zeros.  The connection
formula and the standard sectorial expansion of \(E_1\) give, as
\(x\to+\infty\),
\begin{equation}\label{eq:ssc-H-imaginary-asymptotic}
 H(ix)=\gamma+\log x+\frac{\pi i}{2}+O(x^{-1}),
 \qquad H(-ix)=\overline{H(ix)}.
\end{equation}
The two-exponential formula \eqref{eq:Hprime-sinh}, followed by one
integration, gives the indicator
\[
 h_{E_0}(\vartheta)=\frac12|\sin\vartheta|.
\]
In particular \(E_0\) has exponential type \(1/2\).  Equation
\eqref{eq:ssc-H-imaginary-asymptotic} shows that
\(|E_0(x)|=O(\log(2+|x|))\) on the real axis, so \(E_0\) belongs to the
Cartwright class.

It follows that \(\Theta=E_0^\#/E_0\) is analytic and of bounded type in
the upper half-plane, with unimodular boundary values.  For
\(y\to+\infty\), the real-axis asymptotics obtained from the Joukowski
representation are
\[
 H(y)=\gamma+\log y+O(y^{-1}),\qquad
 H(-y)=-\frac{e^y}{y}\bigl(1+O(y^{-1})\bigr).
\]
Therefore
\begin{equation}\label{eq:ssc-Theta-mean-type}
 \Theta(iy)=e^{-y}\frac{H(-y)}{H(y)}
 =-\frac{1+o(1)}{y\log y}.
\end{equation}
Thus \(\Theta\) has zero mean type.  Bounded-type factorization, together
with its unimodular boundary values, now shows that \(\Theta\) is inner:
meromorphic continuation across the real axis removes the singular-inner
measure, while zero mean type removes the exponential factor in the
canonical factorization \cite{deBranges1968,Clark1972}.  It is nonconstant,
so the Hermite--Biehler inequality is strict.
Meromorphic continuation follows from \eqref{eq:ssc-E0-Theta}.  The
description of its zeros is immediate from \(H(iz)=0\), and the final
assertion is the standard meromorphic-inner Clark/de Branges
correspondence \cite{Clark1972,deBranges1968}.
\end{proof}

\begin{theorem}[Arithmetic independence of the Clark data]
\label{thm:ssc-arithmetic-Clark}
Let \(x_1,\ldots,x_m\in\Qbar\cap\mathbb R^\times\), and suppose that
\(x_j^2\ne x_k^2\) for \(j\ne k\).  Then the \(m\) numbers
\[
 R(x_j):=\frac{H(ix_j)}{H(-ix_j)}
\]
are algebraically independent over \(\Eexp\).  Separately, the \(m\)
numbers
\[
 \Theta(x_j)=e^{ix_j}R(x_j)
\]
are algebraically independent over \(\Eexp\).  In particular each of
\[
 \{R(n):n\geq1\},\qquad \{\Theta(n):n\geq1\}
\]
is a countably algebraically independent family.

If \(\alpha\in\Qbar\cap\mathbb T\), every real solution of
\(\Theta(x)=\alpha\) is transcendental, apart from the unavoidable
solution \(x=0\) when \(\alpha=1\).
\end{theorem}

\begin{proof}
The \(2m\) algebraic points \(ix_j,-ix_j\) are distinct.  Their Joukowski
fibres are pairwise disjoint.  Theorem~\ref{thm:pure-exp-base}, followed by
the same disjoint-support linear projection used in
Theorem~\ref{thm:H-multipoint-descent}, shows that
\[
 H(ix_1),H(-ix_1),\ldots,H(ix_m),H(-ix_m)
\]
are algebraically independent over \(\Eexp\).  Ratios made from disjoint
pairs of independent indeterminates remain algebraically independent:
indeed, a polynomial relation among the ratios becomes, after clearing
the denominator monomial, a nonzero polynomial relation among the original
indeterminates.  Multiplication by the nonzero coefficients
\(e^{ix_j}\in\Eexp\) preserves independence.

If a nonzero algebraic real \(x\) satisfied \(\Theta(x)=\alpha\), then
\[
 e^{ix}H(ix)-\alpha H(-ix)=0
\]
would be a nontrivial linear relation over \(\Eexp\) between two
independent values.  This is impossible.  At \(x=0\), one has
\(\Theta(0)=1\), which explains the unique exception.
\end{proof}

\begin{remark}[Why signs must not be duplicated]
The restriction \(x_j^2\ne x_k^2\) is necessary: on the real axis
\[
 \Theta(-x)=\Theta(x)^{-1},\qquad R(-x)=R(x)^{-1}.
\]
Thus a family containing both \(x\) and \(-x\) has an immediate algebraic
relation.
\end{remark}

\subsection{Recovery of \texorpdfstring{$\gamma$}{gamma} from arithmetic
scattering data}

For \(n\geq1\), put
\begin{equation}\label{eq:ssc-an-Rn-Tn}
 a_n=\frac{n+\sqrt{n^2+4}}2,\qquad
 R_n=R(n)=\frac{H(in)}{H(-in)}=e^{-in}\Theta(n),\qquad
 T_n=\frac{\pi i}{2}\frac{R_n+1}{R_n-1}.
\end{equation}
Theorem~\ref{thm:ssc-arithmetic-Clark} implies \(R_n\ne1\), so \(T_n\)
is well defined.  Since \(|R_n|=1\), it is real.

\begin{theorem}[Scattering recovery and its first correction]
\label{thm:ssc-gamma-scattering}
One has
\begin{equation}\label{eq:ssc-gamma-raw-recovery}
 \gamma=\lim_{n\to\infty}\bigl(T_n-\log n\bigr),
 \qquad
 T_n-\log n
 =\gamma+O\!\left(\frac{\log n}{n}\right).
\end{equation}
More precisely, the completely explicit corrected estimator
\begin{equation}\label{eq:ssc-gamma-corrected-estimator}
 \widehat\gamma_n=
 \left(1-\frac{2(1+\cos a_n)}{\pi a_n}\right)T_n
 +\frac{\sin a_n}{a_n}-\log a_n
\end{equation}
satisfies
\begin{equation}\label{eq:ssc-gamma-corrected-rate}
 \widehat\gamma_n
 =\gamma+O\!\left(\frac{\log n}{n^2}\right).
\end{equation}
Thus \(\gamma\) is recovered from ratios of values generated by
\(\Ein\) at moving algebraic points, even though every finite family of
those ratios is algebraically independent over \(\Eexp\).
\end{theorem}

\begin{proof}
The identity \(a_n-a_n^{-1}=n\) gives \(J(ia_n)=in\), hence
\begin{equation}\label{eq:ssc-Hin-Ein}
 H(in)=\Ein(ia_n)+\Ein(-i/a_n).
\end{equation}
Put \(L_n=\gamma+\log a_n\) and \(b=\pi/2\).  On the positive imaginary
axis, the connection formula and the standard asymptotic expansion
\[
 E_1(ia)=\frac{e^{-ia}}{ia}
 \left(1-\frac1{ia}+O(a^{-2})\right)
\]
combine with
\(\Ein(-i/a)=-i/a+1/(4a^2)+O(a^{-3})\).  Separating real and imaginary
parts in \eqref{eq:ssc-Hin-Ein} gives
\begin{align}
 \Re H(in)
 &=L_n-\frac{\sin a_n}{a_n}
   +\frac{\cos a_n+1/4}{a_n^2}+O(a_n^{-3}),
   \label{eq:ssc-Hin-real-expansion}\\
 \Im H(in)
 &=b-\frac{1+\cos a_n}{a_n}
   -\frac{\sin a_n}{a_n^2}+O(a_n^{-3}).
   \label{eq:ssc-Hin-imag-expansion}
\end{align}
Because \(H(-in)=\overline{H(in)}\), direct algebra gives
\begin{equation}\label{eq:ssc-Tn-real-ratio}
 T_n=b\frac{\Re H(in)}{\Im H(in)}.
\end{equation}
Equations \eqref{eq:ssc-Hin-real-expansion}--\eqref{eq:ssc-Tn-real-ratio}
first yield
\(T_n=L_n+O(L_n/a_n)\).  Since
\(a_n=n+O(n^{-1})\), this proves
\eqref{eq:ssc-gamma-raw-recovery}.

For the sharper statement, multiply \eqref{eq:ssc-Tn-real-ratio} by
\[
 1-\frac{1+\cos a_n}{b a_n}
\]
and add \(\sin(a_n)/a_n\).  The terms of order \(a_n^{-1}\) in both the
numerator and denominator cancel, leaving
\[
 \left(1-\frac{1+\cos a_n}{b a_n}\right)T_n
 +\frac{\sin a_n}{a_n}
 =L_n+O\!\left(\frac{L_n}{a_n^2}\right).
\]
Since \(b=\pi/2\), subtraction of \(\log a_n\) gives
\eqref{eq:ssc-gamma-corrected-estimator} and
\eqref{eq:ssc-gamma-corrected-rate}.
\end{proof}

\subsection{An exact inverse-resonance trace}

\begin{theorem}[Inverse-resonance trace]
\label{thm:ssc-inverse-resonance}
Let \(Z(H)\) denote the multiset of all zeros of \(H\).  Then
\begin{equation}\label{eq:ssc-inverse-resonance}
 \boxed{
 \sum_{\rho\in Z(H)}\frac{-2\Re\rho}{|\rho|^2}
 =1+\frac{\sin1}{\Cin(1)}.}
\end{equation}
The series converges absolutely, every summand is positive, and its value
is transcendental over \(\Eexp\), hence also over
\(\Qbar(e,e^i)\).
\end{theorem}

\begin{proof}
By Theorem~\ref{thm:ssc-HB-Clark}, \(\Theta\) is a meromorphic inner
function of zero mean type.  Meromorphic continuation across the real axis
excludes a finite singular measure in its inner factorization, and zero
mean type excludes an exponential inner factor.  Thus \(\Theta\) is a
pure Blaschke product.  Its zeros are
\[
 \zeta_\rho=-i\rho\in\mathbb C_+
 \qquad(\rho\in Z(H)).
\]
Logarithmic differentiation of the Blaschke product at zero gives
\begin{equation}\label{eq:ssc-Blaschke-log-derivative}
 \frac1i\frac{\Theta'(0)}{\Theta(0)}
 =\sum_{\rho\in Z(H)}
   \frac{2\Im\zeta_\rho}{|\zeta_\rho|^2}
 =\sum_{\rho\in Z(H)}
   \frac{-2\Re\rho}{|\rho|^2}.
\end{equation}
The Blaschke condition gives absolute convergence; equivalently, all terms
are positive and there are no zeros near the origin.

On the other hand, \(\Theta(0)=1\),
\(H(0)=2\Cin(1)\), and \(H'(0)=\sin1\).  Differentiating
\eqref{eq:ssc-E0-Theta} directly yields
\begin{equation}\label{eq:ssc-origin-phase-derivative}
 \frac1i\frac{\Theta'(0)}{\Theta(0)}
 =1+2\frac{H'(0)}{H(0)}
 =1+\frac{\sin1}{\Cin(1)}.
\end{equation}
Together with \eqref{eq:ssc-Blaschke-log-derivative}, this proves the
trace identity.

Finally \(\sin1\in\Eexp^\times\), whereas
\[
 \Cin(1)=\frac{\Ein(i)+\Ein(-i)}2
\]
is transcendental over \(\Eexp\) by
Theorem~\ref{thm:pure-exp-base}.  Hence the right side of
\eqref{eq:ssc-inverse-resonance} is transcendental over \(\Eexp\).
\end{proof}

\begin{remark}[Logical scope]
The stability and Clark results above are unconditional consequences of
the Joukowski descent and the multipoint \(\Ein\) theorem.  The recovery
formula \eqref{eq:ssc-gamma-raw-recovery} is an asymptotic encoding of
\(\gamma\), not an irrationality criterion: passing from algebraic
independence at each finite set of moving points to an arithmetic statement
about the limit requires a genuinely uniform measure.  Likewise
Proposition~\ref{prop:ssc-euler-equivalences} isolates an exact algebraic
witness but does not establish that the witness exists.
\end{remark}
